\section{Primary intervention and retained alternatives}

F3 is the leading \emph{experimental candidate} rather than a confirmed method contribution.
F1 is a possible evaluator and F7 a conditional independent Ising candidate.
The remaining families are deferred, not erased.

The active inventory contains 21 stable idea IDs, consolidated into seven families.
The repository candidate cards give each family a method target, physical case, closest
comparators and falsifier. The following relations define the starting work; they are
not claimed as new mathematics.

\subsection{F1: Event-aware inference-failure estimation}

Represent dominant persistent circuit effects by latent parameters $z$, retain remaining
variation conditionally, partition important directions at actual ADC decoding transitions,
and integrate tractable pieces while sampling the residual. An inference reuses the same
chip realization across all its reads. The difficult contribution is a compact
circuit-derived representation with useful accuracy and scaling.

For signed bit-plane coefficients $a_b$ and read errors $e_b$,
\begin{align}
\Delta y &= \sum_b a_b e_b,\\
\operatorname{Var}(\Delta y)&=\sum_b a_b^2\operatorname{Var}(e_b)
 +2\sum_{b<c}a_ba_c\operatorname{Cov}(e_b,e_c).
\end{align}
Marginal count-error rates do not supply these covariances or determine classification
failure. E06 compares total prediction cost with calibrated strong sampling methods and
compatible verification methods. Boundary explosion, setup overhead or no measurable
benefit is a stopping outcome.

\subsection{F2: Spatially informed polarity}

For 0/1 bit planes, let $K=W^T\mathbf1$. Then
\begin{equation}
 W^Tx=K-W^T(\mathbf1-x).
\end{equation}
The decoding identity must be adapted for XNOR or complementary-cell circuits. If a
validated electrical proxy is $J(a)=a^THa$ with symmetric $H$,
\begin{equation}
 J(x)-J(\mathbf1-x)=(2x-\mathbf1)^TH\mathbf1.
\end{equation}
Precomputing $\lambda=H\mathbf1$ gives a weighted input score. Uniform weights reduce to
ordinary minority inversion. The proxy's link to decoding and total cost remains to be
validated. Column-count metadata alone costs $9\times256=2{,}304$ bits per binary array.

Given input activity $m$ and known column count $K_j$, the possible ideal overlap-count
cardinality is
\begin{equation}
 L=1+\min(m,256-m,K_j,256-K_j),
\end{equation}
which is unchanged by complementation. Lower activity therefore does not automatically
reduce the information that must be resolved. E04 compares calibrated direct reading,
minority inversion, the electrical selector and an offline oracle, including control cost.

\subsection{F3: Reliability-constrained mapping and readout}

A row permutation $P$ applied consistently to $x$ and $W$ preserves ideal output:
\begin{equation}
 (PW)^T(Px)=W^Tx.
\end{equation}
Its electrical effect need not be invariant. For fixed total input activity $M$, $R$ rows
per group and $C$ nonempty groups,
\begin{equation}
 \overline m_{\mathrm{executed}}=M/C,\qquad
 \lceil M/R\rceil\le C\le\min(M,256/R).
\end{equation}
The lower bound permits an input-specific packing oracle, not a universally achievable
static permutation. Increasing mean executed-read activity is a concentration effect;
it does not prove aggregate reliability degradation.

For a simplified expected count-error cost $c(m)$ with $c(0)=0$, combining loads $a,b$
changes the cost by
\begin{equation}
 \Delta c=c(a+b)-c(a)-c(b).
\end{equation}
Linear $c$ leaves the expected total unchanged; a convex cost can penalize concentration.
Real cost depends additionally on pattern, programmed weights, position, calibration and
timing. Task failure further depends on joint errors and signed reconstruction.

An illustrative scalarized read cost $F(m)=E_0\mathbf1_{m>0}+E_1m+\lambda c(m)$ yields
\begin{equation}
 \mathbb E[\Delta F]=-E_0\Pr(a>0,b>0)
 +\lambda\mathbb E[c(a+b)-c(a)-c(b)].
\end{equation}
This separates saved fixed read overhead from a concentration penalty. It is a modelling
framework, not a measured MRAM energy law or an originality claim.

The candidate mapper starts from CAR and balanced orders, scores complete calibration
traces using a validated response/cost model, and accepts bounded row swaps satisfying
reconstructed-error, timing and area budgets. Each baseline may retune its supported
readout. A generic constrained optimizer alone is insufficient; the deliverable must
include a useful circuit-dependent grouping rule and a demonstrated advantage.

\subsubsection{A narrower candidate: covariance of decoded bit-cycle errors}

On 4 November the F3 method question was sharpened: can two mappings with equal read
counts, activities and marginal decoded-error laws have different reconstructed errors
after scalar calibration? For reconstruction coefficients $b$ and decoded error vector
$e$, the proposed score is
\begin{equation}
 R(P)=\mathbb E[(b^Te)^2]=(b^T\mu_e)^2+b^T\Sigma_e b.
\end{equation}
Pair moments suffice for this MSE identity; they do not determine the full error PMF or
classification-failure probability. A validated conditional current model and actual
decoder boundaries would supply the moments. Shared periphery must remain in the model.

The score is a \emph{search surrogate}. A frozen held-out evaluation must separately
report reconstructed arithmetic error, prediction mismatch relative to the clean
quantized network, and ground-truth classification error with paired input/chip uncertainty.
A variance advantage alone does not establish neural reliability. If a column-correlated
estimator is used to rank maps, its omitted joint events and calibration must be audited
against direct persistent-chip sampling.

A single fixed random permutation and the best of random candidates found with the same
offline search budget are \emph{different} baselines. A channel-shuffling comparator is
mandatory, because ordinary decorrelation is already known. Charge legal placement,
routing, capacity, offline optimization, and calibration fairly.

The prospective offline mapper searches bounded row exchanges from CAR and balanced
initializations, updating the affected read moments subject to a fixed resource budget.
One mapping serves every activation plane on the resident array. The general search
implementation remains future work; a two-map score and constructed check are complete.
The distinction from CAR's co-zero affinity \cite{S01}, bit-aware adaptive blocks
\cite{S02}, fault-region mapping \cite{S07} and marginal-code/bit-significance optimization
\cite{S31} is the joint decoded-error score and its grouping consequence. This is a
candidate difference, not a certified absence from prior work. Full DIB-CIM access is
still required.

\subsection{F4: Decision-directed neural refinement}

Begin at the final classifier. A coarse computation supplies score intervals
$[\underline z_j,\overline z_j]$. Accept class $c$ when
\begin{equation}
 \underline z_c>\max_{j\ne c}\overline z_j;
\end{equation}
otherwise refine under a frozen policy. The interval must include omitted terms and
declared hardware error. A statistical envelope requires error control valid under
adaptive refinement. F4 must demonstrate an inexpensive bound and implementation;
bound construction, fallback and the layer's share of total task cost are included.

\subsection{F5: Analog--digital crossover}

Derive the feasible region in which analog cost $E_A(m,R,B,t,z)$ is below digital cost
$E_D(m,z)$ at matched precision, task quality, throughput and capacity. Here $z$ denotes
a declared calibration/chip class. A small mode-selection table is meaningful only if
both modes and their extra periphery are physically supported. E07 compares best fixed
analog, digital and static-hybrid alternatives. A sweep alone is insufficient.

\subsection{F6: Readout identifiability and extra observables}

For one observable $v$, economical readout metadata, class priors $\pi_k$ and calibrated
densities $f_k(v)$, the Bayes error floor is
\begin{equation}
 P_e^*=1-\int\max_k\{\pi_k f_k(v)\}\,dv.
\end{equation}
This is established decision theory. The intended work is a target-specific separation
condition and a minimal additional sample/reference/read configuration that improves
the practical trade-off. Giving the controller the true result, or evaluating the whole
dot product digitally for free, would invalidate that argument. E08 checks identifiability,
noise amplification and calibration alternatives before adding a sensing intervention.

\subsection{F7: Ising field persistence and refinement}

F7a is a conditional diagnostic: distinguish fixed manufacturing variation, physical
temporal read noise, an intentionally decorrelated control and error accumulated in
cached fields. Matching a mean read-error rate does not match conditional error laws.
Full analog field recomputation and incremental reads of binary couplings are different
physical operators. A fixed matrix and consistent initialization make exact increments
telescope; fresh coefficient-read errors can leave a cache residual after a spin flips
and flips back. E09B must distinguish this effect from a temperature-floor hypothesis.
Generic coupling sensitivity and incremental updating are already established
\cite{S18,S32,S33}. A companion chapter remains conditional on a useful specific question,
shared mechanism/method and scope fit; the same noise table alone is insufficient.

Stable binary MRAM can store couplings while randomness is generated separately. For a
monotone update probability $p(h)$, draw one uniform $u$ and obtain a valid field interval
$[L,U]$. Commit the positive update if $u<p(L)$, the negative update if $u\ge p(U)$, and
otherwise refine while retaining $u$. With valid deterministic bounds, the same draw
produces the exact-field decision. Noise, update order, temperature schedule and fallback
must be included. E09 remains a conditional application/companion; refinement's
mathematical comparison and interval hardware remain unresolved.
